\section{Open problems}\label{sec:open}

\begin{enumerate}[label=\arabic*.,leftmargin=2em]
\item \emph{Decide (T).} Either construct a machine transposing $n\times n$ matrices with $\lceil\log_2n\rceil$-bit
entries in time $O(n^2\log n)$, or prove that none exists. A proof is accepted by the frozen checker in the form
described in \cref{sec:model}.

\item \emph{The Tape Coding Conjecture} (\cref{conj:TCC}). A proof refutes~(T) (\cref{thm:reduction}). It suffices to
rule out codes at co-entropy $Cn$ for every $C$, for infinitely many $n$, on diagrams that write their output at the
end (\cref{thm:coent-linear,thm:io}). By \cref{sec:barriers}, a proof must use the shape of tape networks, and its
techniques would give lower bounds for oblivious computation beyond the state of the art. A refutation would be a
family of codes realizing the maximal coding gain for complete exchange on tape networks; it would not by itself give
a machine.

\item \emph{Plain versus robust.} Does the plain conjecture (codes correct on all inputs) imply the robust one? Already
the first rung is open: does ``no code on all inputs'' imply ``no code on half of the inputs''? A positive answer
for codes of trace classes, which are self-detecting, would suffice for~(T) (\cref{sec:plainrobust}).

\item \emph{Oblivious machines.} Is there an oblivious linear-time transposer? The plain conjecture says no
(\cref{thm:oblivious}). We expect that Harvey and van der Hoeven's reduction~\cite{HvdH2025} preserves obliviousness,
so that a negative answer would also rule out oblivious linear-time multiplication; this requires checking that all
control flow and data layout in their construction depend only on the sizes. Does (T) imply oblivious~(T)?

\item \emph{Locally encoded index coding} (\cref{sec:construction}). On bounded-degree expanders, can a broadcast of
two-row functionals of size $(1-\varepsilon)n^2$ let every receiver, knowing its own row, recover its column?

\item \emph{Sharper passes.} Remove the hypothesis $T\le\poly(L)$ from Theorem~P$'$ (for example through the run-length
lemma of Grohe and Schweikardt~\cite{GroheSchweikardt2005}), and extend Theorems~P and~P$'$ to passes that may move
back by a bounded number of cells.

\item \emph{Beyond levels.} Theorem~M* concerns partitions of macroscopic time into intervals. Cuts of other shapes, or
information inequalities that are not cut bounds, are needed to go further; is there a tape-specific non-cut inequality
for oblivious machines?

\item \emph{XOR-frame gadgets.} Is there a gadget (\cref{def:gadget}) with $s/W\le m^{\tau}$ for some $\tau<1$
bounded away from $1$? This would improve the exponent of the sub-$n\log n$ transposition claim substantially, though
not to linear time (\cref{thm:frame-change}). For integer frames the lower bounds of \cref{thm:mod2m} stay below
$4W$, while the trivial gadget costs $mW$; close this gap.
\end{enumerate}
